% Sekvenční diagram pipeline umělé inteligence pro rozpoznání jídla z fotografie.
% Aktéři: UI -- AiPipelineService -- PromptSanitizer -- AiService -- Externí API.
\begin{tikzpicture}[
	font=\scriptsize\sffamily,
	x=1cm, y=1cm,
	actor/.style={
		draw,
		rounded corners=1.5pt,
		minimum height=8mm,
		minimum width=24mm,
		align=center,
		fill=black!8,
	},
	lifeline/.style={dashed, black!50},
	msg/.style={-{Latex[length=2mm]}, thick},
	ret/.style={-{Latex[length=2mm]}, thick, densely dashed},
	mlbl/.style={font=\scriptsize\sffamily, midway, above, sloped=false, fill=white, inner sep=1pt},
	selfnote/.style={
		draw,
		rounded corners=1pt,
		fill=black!4,
		font=\scriptsize\sffamily,
		align=left,
		inner sep=2pt,
	},
]

% --- Aktéři ---
\node[actor] (a1) at (0,0)   {Uživatelské\\rozhraní};
\node[actor] (a2) at (3.4,0) {\texttt{AiPipeline}\\\texttt{Service}};
\node[actor] (a3) at (6.8,0) {\texttt{Prompt}\\\texttt{Sanitizer}};
\node[actor] (a4) at (10.2,0){\texttt{AiService}\\(OpenAI / Gemini)};
\node[actor] (a5) at (13.4,0){Externí API\\(REST)};

% --- Lifelines ---
\foreach \a in {a1,a2,a3,a4,a5}
	\draw[lifeline] (\a.south) -- ($(\a.south)+(0,-10.5)$);

% --- Sekvence zpráv ---
% 1. UI -> Pipeline: vstup
\draw[msg] ($(a1.south)+(0,-0.5)$) -- ($(a2.south)+(0,-0.5)$)
	node[mlbl] {\texttt{analyzeFoodPhoto(photo, text)}};

% 2. Pipeline -> Sanitizer: sanitize(text)
\draw[msg] ($(a2.south)+(0,-1.4)$) -- ($(a3.south)+(0,-1.4)$)
	node[mlbl] {\texttt{sanitize(text)}};

% 3. Sanitizer -> Pipeline: cleanedText (návratová)
\draw[ret] ($(a3.south)+(0,-2.2)$) -- ($(a2.south)+(0,-2.2)$)
	node[mlbl] {\texttt{cleanedText} (regex screening, limit délky)};

% 4. Pipeline self-loop: sestavení JSON promptu
\node[selfnote, anchor=west] at ($(a2.south)+(0.3,-3.2)$) {
	\textbf{sestavení JSON promptu}\\
	(\texttt{lib/utils/prompt.dart}):\\
	\texttt{task} + \texttt{expected\_output} + \texttt{schema.rules}\\
	+ kontext uživatele (dieta, jazyk)
};
\draw[msg] ($(a2.south)+(0,-2.95)$) .. controls +(0.6,0) and +(0.6,0) .. ($(a2.south)+(0,-3.55)$);

% 5. Pipeline -> AiService: generate(prompt + image)
\draw[msg] ($(a2.south)+(0,-4.1)$) -- ($(a4.south)+(0,-4.1)$)
	node[mlbl] {\texttt{generate(prompt, imageBase64)}};

% 6. AiService -> External: HTTPS POST
\draw[msg] ($(a4.south)+(0,-4.9)$) -- ($(a5.south)+(0,-4.9)$)
	node[mlbl] {\texttt{HTTPS POST /chat/completions}};

% 7. External -> AiService: JSON response
\draw[ret] ($(a5.south)+(0,-6.4)$) -- ($(a4.south)+(0,-6.4)$)
	node[mlbl] {JSON odpověď modelu};

% 8. AiService -> Pipeline
\draw[ret] ($(a4.south)+(0,-7.2)$) -- ($(a2.south)+(0,-7.2)$)
	node[mlbl] {\texttt{AiResponse} (deserializace)};

% 9. Pipeline self-loop: confidence gate
\node[selfnote, anchor=west] at ($(a2.south)+(0.3,-8.2)$) {
	\textbf{confidence gate} (práh 0,50):\\
	\(\Rightarrow\) \texttt{success} \(\mid\) \texttt{lowConfidence} \(\mid\) \texttt{failure}
};
\draw[msg] ($(a2.south)+(0,-7.95)$) .. controls +(0.6,0) and +(0.6,0) .. ($(a2.south)+(0,-8.55)$);

% 10. Pipeline -> UI: výsledek
\draw[ret] ($(a2.south)+(0,-9.3)$) -- ($(a1.south)+(0,-9.3)$)
	node[mlbl] {\texttt{AiAnalysisResult}};

% --- Spodní rámečky aktérů (pro ukončení lifeline) ---
\foreach \a/\lbl in {a1/{},a2/{},a3/{},a4/{},a5/{}}
	\node[draw, fill=black!8, minimum height=2mm, minimum width=24mm]
		at ($(\a.south)+(0,-10.5)$) {};

\end{tikzpicture}
